\section{Discussion}
\label{sec:discussion}

\subsection{What this evidence supports}

The evaluation in Section~\ref{sec:evaluation} addresses three separable claims,
and it is worth stating precisely what supports each. \emph{Maturity} is
supported by the test suite, the continuous-integration matrix and the packaging
evidence of Sections~\ref{sec:testing} and~\ref{sec:package-architecture}; the six
defects in Table~\ref{tab:defects} are offered as evidence that testing was
substantive, not that the software is defect-free, which no finite suite can
establish. \emph{Reliability} is supported at two levels that should not be
conflated: behaviour under network failure (Section~\ref{sec:fault-injection}, a
controlled deterministic experiment) and the correctness of the \emph{output},
where the central evidence is the validation study of
Section~\ref{sec:validation-study}. \emph{Value over an unintegrated manual
workflow} is supported by an objective interaction-cost comparison
(Section~\ref{sec:workflow-comparison}), which is a proxy for effort, not a
measurement of human task time or error rate.

\subsection{What validation buys, and what it costs}

The validation study separates two things that a single ``success'' flag
conflates. Accepting the top sequence-search hit returns a structure not
catalogued under the queried accession for most queries in a random sample, but
most of those are proteins that have no deposited structure at all, for which a
homologue comes back; for proteins that do have a structure the naive policy was
wrong for a minority of queries, and a smaller part of those were clearly a
different protein. Validation converts these into explicit levels. The strict
policy accepted no wrong structure, and the permissive policy's false
associations were all the same or a near-identical sequence. The price is
sensitivity: \N{s1-has-sens}\% of proteins with a structure received an M4 against
\N{n1-has-sens}\% for the naive policy, because many correct structures are partial,
cannot be validated, or lack a confirmed accession.

We therefore present the two operating points as a choice, not a verdict. A user
who needs to assert that a structure \emph{is} the queried protein should take M4
and accept that it will often return nothing; a user who can use the same or a
near-identical sequence should take M3 or better; a user who simply wants a
related fold can take any level but should read the level. Partial structures,
now flagged, are real structures of the protein and are often what a
structural biologist wants; the flag lets them be found without weakening M4.

\subsection{The AlphaFold defect as a general lesson}

The AlphaFold DB retrieval defect (Table~\ref{tab:defects};
Section~\ref{sec:live-benchmark}) is discussed not because it was unusually severe
but because of how it presented: a client that built the model URL from an assumed
version number returned HTTP 404 for every request once the database was
republished, which is indistinguishable in isolation from the absence of a model.
The general lesson is that a retrieval client's assumptions about a remote
resource should be verified against its authoritative interface, and that a
failure rate near 100\% on a resource that should have high coverage is itself
a signal. The same applies to this study's own tooling, as the next subsection
records.

\subsection{Verifying the verifier}

The validation study itself went wrong in ways that are instructive, and we
report them because the same errors are easy to make in any benchmark of a
validation layer. Our first analysis scored the naive baselines as accepted
exactly when their hit was correct and then evaluated them against that same
label, which made them perfect by construction. Our first threshold analysis
defined its ``true positives'' as the classifier's own M4 output, which made the
separation circular. An early sampler took the first results returned by a search
and produced an all-human sample. The adversarial set was built with an ordering
that depended on string hashing, so a ``seeded'' draw differed between runs. A
documentation table stated thresholds that did not match the shipped defaults.
Each was found by inspecting the output against what it should have been, was
corrected before the numbers reported here were produced, and is logged as a
protocol amendment in the repository. The practical conclusion is that a
benchmark of a verification step needs the same adversarial scrutiny as the step.

\subsection{Why the classification generalises, and where it does not}

The M0--M4 design is deliberately resource-agnostic: it depends only on a query
sequence, a returned chain, and optionally a set of cross-references, so it applies
unchanged to experimental structures, accession-keyed predictions and
sequence-only predictions. Its strongest levels depend on cross-references that
only some resources provide, which is why ESM Atlas predictions cannot reach M4 by
design. Its weakest level, M2, depends on an identity measure that is not
discriminating at the threshold we ship (Section~\ref{sec:threshold}); a user
should treat M2 as a prompt to look, not as a finding.
